\section{Introduction} \label{sec:introduction}
Foams are among the most familiar yet physically complex systems encountered in everyday life. From the froth on your morning coffee to the structure of biological tissues, foams occupy a unique place in the study of soft matter. A foam is a layer of gas bubbles seperated by a liquid medium, and the collective behaviour of these bubbles; how they coarsen over time and how the foam restructures because of this, is the core of what is being investigated in this report.\\

Coarsening is a fundamental property of these foams. It is the process by which bubbles exchange gas across their shared interfaces causing larger bubbles to get larger and smaller bubbles get smaller. Over time, the average bubble area increases and the number of bubbles decrease. During this coarsening, once the small bubbles disappear, the larger gain new neighbours and the foam experiences a restructure. Understanding the statistical character of the individual bubble trajectories, and they type of randomness they follow, within this evolving foam structure is the main motivation of this project.\\

In this project, we use data from a 2D foam simulated by Prof. Simon Cox. In this data we are given the positions of the bubble centres at each timestep, along with individual properties such as area and coordination number. Tools from econophysics are borrowed to investigate the dynamics and overall statisticss of this foam. The central analogy is between the bubbles' displacements through the foam and the price movements in a financial market. In both cases, one is interested in the nature of the randomness involved. Whether it is completely random with simple Brownian diffusion, and whether there is a need to characterise a different type of randomness.\\

Brownian motion is used in this project as a baseline of randomness for comparison. It is well understood in it's statistical properties; Gaussian displacement distributions, mean squared displacement that scales linearly in time, and vanishing temporal autocorrelation. This means we can use it as a sanity check for all of our methods, if we apply our methods to it and get the expected results, we can be more confident in our results when we apply the same methods to the bubble trajectories. Deviations from this baseline are of particular interest as they suggest the presence of underlying statistical properties of the foam that influence individual bubble motion.\\

The report is structured as follows: Section \ref{sec:theoreticalbackground} reviews the statistical mechanics of Brownian motion, the tools from econophysics used in this projec; log-returns, volatility, temporal autocorrelation, etc., and how they can be intuitively applied to the bubble displacements. Section \ref{sec:compmethod} describes the use of pandas to organise the data in a data-frame, extract information, and how computational methods are applied to this information. Section \ref{sec:resultsanddiscussion}, presents the results and discuss their implications for understanding bubble motion in coarsening foams. Finally, section \ref{sec:conclusion} summarises our findings and suggests directions for future research.
